\begin{center}
\Large\bfseries Appendices
\end{center}
\vspace{1em}

\section{Dominick's Dataset Details}
\label{app:dominicks}

The Dominick's Finer Foods dataset \citep{dominicks1999} is maintained by the James M. Kilts Center for Marketing at the University of Chicago Booth School of Business. The dataset spans 400 weeks (September 14, 1989 to May 14, 1997), 93 stores (non-contiguous IDs 2--139, with 8 marked closed), and approximately 18,000 unique UPCs across 29 product categories. The movement files record store-SKU-week tuples with unit sales (\texttt{MOVE}), prices (\texttt{PRICE}), profit margins (\texttt{PROFIT}), and promotion indicators (\texttt{SALE}). The store demographic file (\texttt{demo.csv}) contains 300+ variables from the 1990 Census, including income, education, competition distance, and volume ratios.

Table~\ref{tab:category-codes} lists the 29 product category codes and their descriptions. The canned soup category (\texttt{cso}) serves as the primary evaluation target in this work. Other recommended categories for future work include beer (\texttt{ber}), soft drinks (\texttt{sdr}), and cereals (\texttt{cer}).

\begin{table}[h]
\centering
\caption{Dominick's product category codes and descriptions.}
\label{tab:category-codes}
\begin{tabular}{@{}lll@{}}
\toprule
\textbf{Code} & \textbf{Description} & \textbf{Notes} \\
\midrule
\texttt{ana} & Analgesics & \\
\texttt{bat} & Bath products & \\
\texttt{ber} & Beer & High cross-elasticity, $\sim$791 UPCs \\
\texttt{bjc} & Bottled juice & \\
\texttt{cer} & Cereals & \\
\texttt{che} & Cheese & \\
\texttt{cig} & Cigarettes & \\
\texttt{coo} & Cookies & \\
\texttt{cra} & Crackers & \\
\texttt{cso} & Canned soup & Primary evaluation; $\sim$25 SKUs \\
\texttt{did} & Dish detergent & \\
\texttt{fec} & Fabric conditioner & \\
\texttt{frd} & Frozen dinners & \\
\texttt{fre} & Frozen entrees & \\
\texttt{frj} & Frozen juice & \\
\texttt{fsf} & Front-end snacks & \\
\texttt{gro} & Grooming & \\
\texttt{lnd} & Laundry detergent & \\
\texttt{oat} & Oatmeal & \\
\texttt{ptw} & Paper towels & \\
\texttt{sdr} & Soft drinks & High promotional activity \\
\texttt{sha} & Shampoo & \\
\texttt{sna} & Snacks & \\
\texttt{soa} & Soap & \\
\texttt{tbr} & Toothbrush & \\
\texttt{tna} & Canned tuna & \\
\texttt{tpa} & Toothpaste & \\
\texttt{tti} & Toilet tissue & \\
\bottomrule
\end{tabular}
\end{table}

Table~\ref{tab:data-statistics} summarizes key statistics for the canned soup category after preprocessing. Rows with \texttt{OK} = 0 or \texttt{PRICE} $\leq$ 0 are dropped. \texttt{PRICE\_HEX} and \texttt{PROFIT\_HEX} columns are discarded. Unit price is computed as \texttt{PRICE} / \texttt{QTY}; wholesale cost is \texttt{PRICE} $\times$ (1 $-$ \texttt{PROFIT}/100) / \texttt{QTY}. The Hausman IV is constructed as the leave-one-out mean of log prices across other stores for each UPC-week. The temporal split is strictly chronological: train weeks 1--280, validation 281--340, test 341--400.

\begin{table}[h]
\centering
\caption{Data statistics for canned soup category (preprocessed).}
\label{tab:data-statistics}
\begin{tabular}{@{}lr@{}}
\toprule
\textbf{Statistic} & \textbf{Value} \\
\midrule
Training tuples & $\sim$581,000 \\
Validation tuples & $\sim$73,000 \\
Test tuples & $\sim$124,500 \\
Unique UPCs & $\sim$25 \\
Stores (after filtering) & 83 \\
Train weeks & 1--280 \\
Validation weeks & 281--340 \\
Test weeks & 341--400 \\
\bottomrule
\end{tabular}
\end{table}

\section{Hyperparameter Configuration}
\label{app:hyperparams}

Table~\ref{tab:hyperparams-full} provides the complete hyperparameter configuration for \dreamprice{}. World model parameters include the \mamba{} backbone dimension (\texttt{d\_model}), stochastic latent specification (\texttt{n\_cat}, \texttt{n\_cls}), and reward ensemble size (\texttt{n\_heads}). Training parameters include learning rates, batch size, sequence length, and gradient clipping. Reinforcement learning parameters include the discount factor $\gamma$, lambda-return $\lambda$, entropy scale $\eta$, MOPO LCB coefficient $\lambda_{\text{LCB}}$, and imagination horizon $H$. Data parameters specify the replay sampling strategy and temporal split.

\begin{table}[h]
\centering
\caption{Full hyperparameter configuration.}
\label{tab:hyperparams-full}
\small
\begin{tabular}{@{}lll@{}}
\toprule
\textbf{Parameter} & \textbf{Value} & \textbf{Notes} \\
\midrule
\multicolumn{3}{l}{\textit{World model}} \\
\texttt{d\_model} & 512 & Mamba-2 backbone dimension \\
\texttt{n\_cat} & 32 & Categorical distributions \\
\texttt{n\_cls} & 32 & Classes per categorical \\
\texttt{n\_heads} & 5 & Reward ensemble for MOPO LCB \\
\texttt{unimix} & 0.01 & Uniform mixing on categoricals \\
\texttt{twohot\_bins} & 255 & Symlog range $[-20, +20]$ \\
\midrule
\multicolumn{3}{l}{\textit{Training}} \\
\texttt{lr\_world\_model} & $1 \times 10^{-4}$ & Adam \\
\texttt{lr\_actor} & $3 \times 10^{-5}$ & Adam \\
\texttt{lr\_critic} & $3 \times 10^{-5}$ & Adam \\
\texttt{batch\_size} & 32 & Sequences per step \\
\texttt{seq\_len} & 64 & Time steps per sequence \\
\texttt{grad\_clip\_wm} & 1000 & Global norm \\
\texttt{grad\_clip\_ac} & 100 & Global norm \\
\texttt{$\beta_{\text{pred}}$} & 1.0 & Prediction loss weight \\
\texttt{$\beta_{\text{dyn}}$} & 0.5 & Dynamics loss (sg posterior) \\
\texttt{$\beta_{\text{rep}}$} & 0.1 & Representation loss (sg prior) \\
\texttt{free\_bits} & 1.0 & Nat minimum KL \\
\midrule
\multicolumn{3}{l}{\textit{Reinforcement learning}} \\
$\gamma$ & 0.95 & Discount factor \\
$\lambda$ & 0.95 & Lambda-return \\
$\eta$ & $3 \times 10^{-4}$ & Entropy regularization \\
$\lambda_{\text{LCB}}$ & 1.0 & MOPO pessimism \\
$H$ & 13 & Imagination horizon (steps) \\
\texttt{critic\_ema\_decay} & 0.98 & Slow critic \\
\texttt{return\_norm\_ema} & 0.99 & Percentile normalization \\
\midrule
\multicolumn{3}{l}{\textit{Data}} \\
\texttt{replay\_uniform} & 70\% & Across quarterly strata \\
\texttt{replay\_recent} & 30\% & Overweight last 2 years \\
\texttt{train\_weeks} & 1--280 & \\
\texttt{val\_weeks} & 281--340 & \\
\texttt{test\_weeks} & 341--400 & \\
\bottomrule
\end{tabular}
\end{table}

\section{Additional Ablation Results}
\label{app:ablation-horizon}

Table~\ref{tab:horizon-sweep} provides the detailed horizon sweep results. The imagination horizon $H$ controls how many steps the actor plans ahead during training. The default $H = 13$ aligns with the retail quarter (4-5-4 calendar). Each configuration uses seed $s = 42$ over 100K gradient steps.

\begin{table}[h]
\centering
\caption{Imagination horizon sweep. Mean return and WM loss for cumulative gross margin over test period. Single seed ($s = 42$).}
\label{tab:horizon-sweep}
\begin{tabular}{@{}lccl@{}}
\toprule
\textbf{Horizon $H$} & \textbf{Return} & \textbf{WM Loss} & \textbf{Notes} \\
\midrule
5 & 145.6 & 22.51 & Short-term exploitation \\
10 & 86.4 & 22.43 & Intermediate transition \\
13 (default) & \textbf{193.7} & 22.44 & Retail quarter alignment \\
25 & 135.3 & 22.48 & Long-horizon planning \\
\bottomrule
\end{tabular}
\end{table}

\section{Computational Requirements}
\label{app:compute}

Table~\ref{tab:compute} summarizes observed training times per seed. All runs use the NVIDIA DGX Spark (GB10 GPU, 128\,GB unified memory). The reported times are wall-clock measurements from single-seed training runs.

\begin{table}[h]
\centering
\caption{Observed training time per seed (hours). Hardware: NVIDIA DGX Spark, 128\,GB unified memory. Ablation times measured from completed runs.}
\label{tab:compute}
\begin{tabular}{@{}lc@{}}
\toprule
\textbf{Configuration} & \textbf{Hours/seed} \\
\midrule
\dreamprice{} (main, 100K steps) & 2.6 \\
\midrule
\multicolumn{2}{l}{\textit{Ablations (single seed, 100K steps)}} \\
Imagination OFF & 0.9 \\
No MOPO LCB & 4.4 \\
Deterministic-only latent & 4.1 \\
No symlog+twohot & 4.4 \\
Horizon sweep ($H \in \{5, 10, 25\}$) & 2.5--4.4 \\
GRU backbone & 2.0 \\
Flat encoder & 2.5 \\
\bottomrule
\end{tabular}
\end{table}

\section{RL Baseline Hyperparameters}
\label{app:rl-baselines}

Table~\ref{tab:rl-baseline-hyperparams} lists the hyperparameters used for the model-free RL baselines. All baselines are implemented via Stable-Baselines3 \citep{raffin2021sb3} and trained within the learned \dreamprice{} world model for 50{,}000 timesteps with a single seed ($s = 42$).

\begin{table}[h]
\centering
\caption{Model-free RL baseline hyperparameters (Stable-Baselines3).}
\label{tab:rl-baseline-hyperparams}
\small
\begin{tabular}{@{}llll@{}}
\toprule
\textbf{Parameter} & \textbf{DQN} & \textbf{PPO} & \textbf{SAC} \\
\midrule
Timesteps & 50{,}000 & 50{,}000 & 50{,}000 \\
Learning rate & $1 \times 10^{-3}$ & $3 \times 10^{-4}$ & $3 \times 10^{-4}$ \\
Batch size & 64 & 64 & 256 \\
Buffer size & 10{,}000 & --- & 10{,}000 \\
Exploration fraction & 0.3 & --- & --- \\
Target update interval & 500 & --- & --- \\
\texttt{n\_steps} & --- & 256 & --- \\
\texttt{n\_epochs} & --- & 10 & --- \\
Clip range & --- & 0.2 & --- \\
$\tau$ (soft update) & --- & --- & 0.005 \\
Entropy tuning & --- & --- & auto \\
Action space & discrete (21 levels) & continuous & continuous \\
\bottomrule
\end{tabular}
\end{table}

Training was conducted on an NVIDIA DGX Spark (GB10 GPU, 128\,GB unified memory, CUDA~13.0, ARM64) using Docker containers for reproducibility. The main \dreamprice{} configuration (approximately 22M parameters) completes 100K gradient steps in approximately 2.6 hours wall-clock time with batch size 32 and sequence length 64. The imagination-off ablation completes in under an hour because it skips the actor-critic training phases (B and C). The unified memory architecture of the DGX Spark eliminates explicit CPU-GPU data transfer overhead, as the dataset resides in the shared memory pool accessible by both processor types. Memory management required early store filtering and demographic column pruning during data loading to keep peak memory usage below 100\,GB, leaving headroom for the CUDA runtime, Triton JIT compilation, and experiment logging.
